% ============================================================================
% EMCH 792: Optimal State Estimation - Homework 1 (Enhanced Solutions)
% ============================================================================

\documentclass[12pt]{article}

% ============================================================================
% PACKAGE IMPORTS
% ============================================================================
\usepackage{amsmath}
\usepackage{graphicx}
\usepackage{amssymb}
\usepackage{tikz}
\usepackage[margin=1in]{geometry}
\usepackage{setspace}
\usepackage{xcolor}
\usepackage{enumitem}
\usepackage{tcolorbox}
\usepackage{fancyhdr}
\usepackage{lastpage}
\usepackage{booktabs}
\usepackage{array}
\usepackage{colortbl}
\usepackage{bm}

% ============================================================================
% HEADER AND FOOTER SETUP
% ============================================================================
\pagestyle{fancy}
\setlength{\headheight}{14.5pt}
\fancyhf{}

\fancyhead[L]{Page \thepage\ of \pageref{LastPage}}
\fancyhead[C]{EMCH 792}
\fancyhead[R]{JC Vaught}

\renewcommand{\headrulewidth}{0pt}

% ============================================================================
% USC BRAND COLORS
% ============================================================================
\definecolor{USC_Garnet}{HTML}{73000A}
\definecolor{USC_Sandstorm}{HTML}{FFF2E3}
\definecolor{USC_Black90}{HTML}{565656}
\definecolor{USC_Black10}{HTML}{ECECEC}
\definecolor{USC_Honeycomb}{HTML}{65780B}
\definecolor{USC_Atlantic}{HTML}{466A9F}
\definecolor{USC_Congaree}{HTML}{1F414D}
\definecolor{USC_Rose}{HTML}{CC2E40}

% ============================================================================
% BOLD COLORED MATH COMMANDS
% ============================================================================
\newcommand{\Ag}{{\color{USC_Garnet}\bm{A}}}
\newcommand{\Bb}{{\color{USC_Atlantic}\bm{B}}}
\newcommand{\Cc}{{\color{USC_Congaree}\bm{C}}}
\newcommand{\Iv}{{\color{USC_Honeycomb}\bm{I}}}
\newcommand{\vv}{{\color{USC_Honeycomb}\bm{v}}}
\newcommand{\Tr}{{\color{USC_Rose}\textbf{Tr}}}

% Colored matrices (no bold for environments)
\newcommand{\matg}[1]{{\color{USC_Garnet}#1}}
\newcommand{\matb}[1]{{\color{USC_Atlantic}#1}}
\newcommand{\matc}[1]{{\color{USC_Congaree}#1}}
\newcommand{\mH}[1]{{\color{USC_Honeycomb}#1}}
\newcommand{\matr}[1]{{\color{USC_Rose}#1}}

% ============================================================================
% CUSTOM COLORED BOXES
% ============================================================================
\newtcolorbox{problembox}{
  colback=USC_Sandstorm,
  colframe=USC_Garnet,
  fonttitle=\bfseries,
  title=Problem Statement,
  sharp corners,
  colbacktitle=USC_Garnet,
  coltitle=white,
}

\newtcolorbox{stepbox}{
  colback=white,
  colframe=USC_Garnet,
  fonttitle=\bfseries,
  title=Step,
  sharp corners,
  colbacktitle=USC_Garnet,
  coltitle=white,
}

\newtcolorbox{resultsbox}{
  colback=white,
  colframe=USC_Honeycomb,
  fonttitle=\bfseries,
  title=Results,
  sharp corners,
  colbacktitle=USC_Honeycomb,
  coltitle=white,
}

% ============================================================================
% CUSTOM QUESTION COMMAND
% ============================================================================
\newcommand{\question}[1]{%
  \clearpage
  \vspace{0.5cm}
  {\noindent\normalsize \textbf{#1}}
  \vspace{0.2cm}
  \hrule
  \vspace{0.1cm}
  \hrule
  \vspace{0.3cm}
}

% ============================================================================
% DOCUMENT INFORMATION
% ============================================================================

\title{EMCH 792: Optimal State Estimation \\ Homework 1 \\ \large{Written Portion Solutions}}
\author{Instructor: Dr.\ Yimin Zhang, Department of Mechanical Engineering \\ University of South Carolina \\ \textit{Solutions by: JC Vaught}}
\date{Due: January 2026}

% ============================================================================
% BEGIN DOCUMENT
% ============================================================================

\begin{document}

\maketitle
\setlength{\parindent}{0pt}

% ============================================================================
% TABLE OF CONTENTS
% ============================================================================

\begin{center}
\begin{tcolorbox}[colback=white, colframe=gray!50!black, colbacktitle=gray!20!white,
                  coltitle=black, sharp corners, boxrule=1pt,
                  title=\Large\bfseries Table of Contents]
\begin{tabular}{p{0.82\textwidth}r}
\textbf{Problem 1.2: Commuting Matrices} \dotfill & \textbf{\pageref{quest:1.2}} \\[0.3em]
\textbf{Problem 1.3: Matrix Identities} \dotfill & \textbf{\pageref{quest:1.3}} \\[0.3em]
\textbf{Problem 1.4: Trace Derivative} \dotfill & \textbf{\pageref{quest:1.4}} \\[0.3em]
\textbf{Problem 1.6: Block Matrix Equation} \dotfill & \textbf{\pageref{quest:1.6}} \\[0.3em]
\textbf{Problem 1.7: Symmetric Matrix Product} \dotfill & \textbf{\pageref{quest:1.7}} \\
\end{tabular}
\vspace{0.2cm}
\end{tcolorbox}
\end{center}

\vspace{0.5cm}

% ============================================================================
% PROBLEM 1.2
% ============================================================================

\question{Problem 1.2: Commuting Matrices}\label{quest:1.2}

\begin{problembox}
Find two $2 \times 2$ matrices $A$ and $B$ such that $A \neq B$, neither $A$ nor $B$ are diagonal, $A \neq cB$ for any scalar $c$, and $AB = BA$. Find the eigenvectors of $A$ and $B$. Note that they share an eigenvector. Interestingly, every pair of commuting matrices shares at least one eigenvector.
\end{problembox}

\begin{stepbox}
\textbf{Part (a): Finding Commuting Matrices}

We seek two $2 \times 2$ matrices $\Ag$ and $\Bb$ such that $\Ag \neq \Bb$, neither is diagonal, $\Ag \neq c\Bb$ for any scalar $c$, and $\Ag\Bb = \Bb\Ag$.

Consider upper triangular matrices with the same diagonal entries:
\begin{equation}
\Ag = \matg{\begin{bmatrix} 1 & 1 \\ 0 & 1 \end{bmatrix}}, \quad \Bb = \matb{\begin{bmatrix} 1 & 2 \\ 0 & 1 \end{bmatrix}}
\end{equation}
\end{stepbox}

\begin{stepbox}
\textbf{Verification of Commutativity}

Compute $\Ag\Bb$:
\begin{equation}
\Ag\Bb = \matg{\begin{bmatrix} 1 & 1 \\ 0 & 1 \end{bmatrix}} \matb{\begin{bmatrix} 1 & 2 \\ 0 & 1 \end{bmatrix}} = \matc{\begin{bmatrix} 1 & 3 \\ 0 & 1 \end{bmatrix}}
\end{equation}

Compute $\Bb\Ag$:
\begin{equation}
\Bb\Ag = \matb{\begin{bmatrix} 1 & 2 \\ 0 & 1 \end{bmatrix}} \matg{\begin{bmatrix} 1 & 1 \\ 0 & 1 \end{bmatrix}} = \matc{\begin{bmatrix} 1 & 3 \\ 0 & 1 \end{bmatrix}}
\end{equation}

Therefore, $\Ag\Bb = \Bb\Ag$. Also: $\Ag \neq \Bb$, neither is diagonal, and $\Ag \neq c\Bb$ for any scalar $c$.
\end{stepbox}

\begin{stepbox}
\textbf{Part (b): Finding Eigenvectors}

\textbf{For matrix $\Ag$:} Solve $(\Ag - \lambda \Iv)\vv = \mH{0}$:
\begin{equation}
\det(\Ag - \lambda \Iv) = \det \begin{bmatrix} 1-\lambda & 1 \\ 0 & 1-\lambda \end{bmatrix} = \matg{(1-\lambda)^2 = 0}
\end{equation}

Thus, $\matg{\lambda = 1}$ is a repeated eigenvalue. For $\lambda = 1$:
\begin{equation}
(\Ag - \Iv)\vv = \begin{bmatrix} 0 & 1 \\ 0 & 0 \end{bmatrix} \begin{bmatrix} v_1 \\ v_2 \end{bmatrix} = \mH{\begin{bmatrix} 0 \\ 0 \end{bmatrix}}
\end{equation}

This gives $v_2 = 0$, so the eigenvector is $\vv_{\Ag} = \mH{\begin{bmatrix} 1 \\ 0 \end{bmatrix}}$.

\textbf{For matrix $\Bb$:} Similarly, $\matb{\lambda = 1}$ is a repeated eigenvalue, and $\vv_{\Bb} = \mH{\begin{bmatrix} 1 \\ 0 \end{bmatrix}}$.
\end{stepbox}

\begin{resultsbox}
\textbf{Final Answers:}

\textbf{Commuting Matrices:}
\[
\Ag = \matg{\begin{bmatrix} 1 & 1 \\ 0 & 1 \end{bmatrix}}, \quad \Bb = \matb{\begin{bmatrix} 1 & 2 \\ 0 & 1 \end{bmatrix}}
\]

\textbf{Shared Eigenvector:}
\[
\boxed{\vv = \mH{\begin{bmatrix} 1 \\ 0 \end{bmatrix}}}
\]

This demonstrates the fundamental result that commuting matrices share at least one eigenvector.
\end{resultsbox}


% ============================================================================
% PROBLEM 1.3
% ============================================================================

\question{Problem 1.3: Matrix Identities}\label{quest:1.3}

\begin{problembox}
Prove the three identities of Equation (1.26):
\begin{enumerate}
    \item $(AB)^T = B^T A^T$
    \item $(AB)^{-1} = B^{-1} A^{-1}$
    \item $\text{Tr}(AB) = \text{Tr}(BA)$
\end{enumerate}
\end{problembox}

\begin{stepbox}
\textbf{Identity 1: $\matg{(\Ag\Bb)^T} = \Bb^T \Ag^T$}

Let $\Ag$ be an $m \times n$ matrix and $\Bb$ be an $n \times p$ matrix. Let $\Cc = \Ag\Bb$.

By definition of matrix multiplication and transpose:
\begin{equation}
\matc{C_{ij}} = \sum_{k=1}^{n} \matg{A_{ik}} \matb{B_{kj}}
\end{equation}

\begin{equation}
(\Ag\Bb)^T_{ij} = \matc{C_{ji}} = \sum_{k=1}^{n} \matg{A_{jk}} \matb{B_{ki}}
\end{equation}

For $\Bb^T \Ag^T$:
\begin{equation}
(\Bb^T \Ag^T)_{ij} = \sum_{k=1}^{n} \matb{B^T_{ik}} \matg{A^T_{kj}} = \sum_{k=1}^{n} \matb{B_{ki}} \matg{A_{jk}} = \sum_{k=1}^{n} \matg{A_{jk}} \matb{B_{ki}}
\end{equation}

Comparing: $(\Ag\Bb)^T_{ij} = (\Bb^T \Ag^T)_{ij}$ for all $i,j$. Therefore: $\boxed{\matg{(\Ag\Bb)^T = \Bb^T \Ag^T}}$
\end{stepbox}

\begin{stepbox}
\textbf{Identity 2: $\matg{(\Ag\Bb)^{-1}} = \Bb^{-1} \Ag^{-1}$}

Assume $\Ag$ and $\Bb$ are invertible square matrices. We verify that $\Bb^{-1} \Ag^{-1}$ is the inverse of $\Ag\Bb$:

\begin{equation}
(\Ag\Bb)(\Bb^{-1} \Ag^{-1}) = \Ag(\Bb\Bb^{-1})\Ag^{-1} = \Ag \Iv \Ag^{-1} = \Ag\Ag^{-1} = \Iv
\end{equation}

\begin{equation}
(\Bb^{-1} \Ag^{-1})(\Ag\Bb) = \Bb^{-1}(\Ag^{-1}\Ag)\Bb = \Bb^{-1} \Iv \Bb = \Bb^{-1}\Bb = \Iv
\end{equation}

Since both products equal the identity: $\boxed{\matg{(\Ag\Bb)^{-1} = \Bb^{-1} \Ag^{-1}}}$
\end{stepbox}

\begin{stepbox}
\textbf{Identity 3: $\Tr(\Ag\Bb) = \Tr(\Bb\Ag)$}

Let $\Ag$ be $m \times n$ and $\Bb$ be $n \times m$.

\begin{equation}
\Tr(\Ag\Bb) = \sum_{i=1}^{m} (\Ag\Bb)_{ii} = \sum_{i=1}^{m} \sum_{j=1}^{n} \matg{A_{ij}} \matb{B_{ji}}
\end{equation}

\begin{equation}
\Tr(\Bb\Ag) = \sum_{j=1}^{n} (\Bb\Ag)_{jj} = \sum_{j=1}^{n} \sum_{i=1}^{m} \matb{B_{ji}} \matg{A_{ij}}
\end{equation}

Since scalar multiplication is commutative and summation order can be exchanged:
\begin{equation}
\Tr(\Bb\Ag) = \sum_{i=1}^{m} \sum_{j=1}^{n} \matg{A_{ij}} \matb{B_{ji}} = \Tr(\Ag\Bb)
\end{equation}

Therefore: $\boxed{\matr{\Tr(\Ag\Bb) = \Tr(\Bb\Ag)}}$
\end{stepbox}

\begin{resultsbox}
\textbf{Summary of Proven Identities:}
\begin{enumerate}
    \item $\boxed{(AB)^T = B^T A^T}$ --- Transpose of a product reverses order
    \item $\boxed{(AB)^{-1} = B^{-1} A^{-1}}$ --- Inverse of a product reverses order
    \item $\boxed{\text{Tr}(AB) = \text{Tr}(BA)}$ --- Trace is invariant under cyclic permutations
\end{enumerate}
\end{resultsbox}


% ============================================================================
% PROBLEM 1.4
% ============================================================================

\question{Problem 1.4: Trace Derivative}\label{quest:1.4}

\begin{problembox}
Find the partial derivative of the trace of $AB$ with respect to $A$.
\end{problembox}

\begin{stepbox}
\textbf{Setup}

Let $\Ag$ be $m \times n$ and $\Bb$ be $n \times m$. The trace of $\Ag\Bb$ is:
\begin{equation}
\Tr(\Ag\Bb) = \sum_{i=1}^{m} (\Ag\Bb)_{ii} = \sum_{i=1}^{m} \sum_{j=1}^{n} \matg{A_{ij}} \matb{B_{ji}}
\end{equation}
\end{stepbox}

\begin{stepbox}
\textbf{Computing the Derivative}

To find the derivative with respect to $\Ag$, compute the derivative with respect to each element $\matg{A_{kl}}$:
\begin{equation}
\frac{\partial \Tr(\Ag\Bb)}{\partial \matg{A_{kl}}} = \frac{\partial}{\partial \matg{A_{kl}}} \sum_{i=1}^{m} \sum_{j=1}^{n} \matg{A_{ij}} \matb{B_{ji}}
\end{equation}

Since only the term with $i=k$ and $j=l$ contains $\matg{A_{kl}}$:
\begin{equation}
\frac{\partial \Tr(\Ag\Bb)}{\partial \matg{A_{kl}}} = \matb{B_{lk}}
\end{equation}

This is precisely the $(k,l)$ element of $\Bb^T$.
\end{stepbox}

\begin{resultsbox}
\textbf{Final Answer:}
\[
\boxed{\frac{\partial \Tr(\Ag\Bb)}{\partial \Ag} = \Bb^T}
\]

The gradient of $\Tr(\Ag\Bb)$ with respect to $\Ag$ is the transpose of $\Bb$.
\end{resultsbox}


% ============================================================================
% PROBLEM 1.6
% ============================================================================

\question{Problem 1.6: Block Matrix Equation}\label{quest:1.6}

\begin{problembox}
Suppose that $A$ is invertible and
\[
\begin{bmatrix} A & A \\ B & A \end{bmatrix} \begin{bmatrix} A \\ C \end{bmatrix} = \begin{bmatrix} 0 \\ I \end{bmatrix}
\]
Find $B$ and $C$ in terms of $A$.
\end{problembox}

\begin{stepbox}
\textbf{Expanding the Block Matrix Multiplication}

\textbf{First row:}
\[
\Ag \cdot \Ag + \Ag \cdot \Cc = \mH{0} \quad \Rightarrow \quad \Ag^2 + \Ag\Cc = \mH{0}
\]

\textbf{Second row:}
\[
\Bb \cdot \Ag + \Ag \cdot \Cc = \Iv \quad \Rightarrow \quad \Bb\Ag + \Ag\Cc = \Iv
\]
\end{stepbox}

\begin{stepbox}
\textbf{Solving for $\Cc$}

From the first row: $\Ag\Cc = -\Ag^2$

Since $\Ag$ is invertible, multiply both sides by $\Ag^{-1}$ on the left:
\[
\Cc = -\Ag
\]
\end{stepbox}

\begin{stepbox}
\textbf{Solving for $\Bb$}

Substituting $\Cc = -\Ag$ into the second row:
\[
\Bb\Ag + \Ag(-\Ag) = \Iv \quad \Rightarrow \quad \Bb\Ag - \Ag^2 = \Iv \quad \Rightarrow \quad \Bb\Ag = \Iv + \Ag^2
\]

Multiplying both sides by $\Ag^{-1}$ on the right:
\[
\Bb = (\Iv + \Ag^2)\Ag^{-1} = \Ag^{-1} + \Ag
\]
\end{stepbox}

\begin{stepbox}
\textbf{Verification}

Check: $\Bb\Ag = (\Ag^{-1} + \Ag)\Ag = \Ag^{-1}\Ag + \Ag^2 = \Iv + \Ag^2$ \checkmark
\end{stepbox}

\begin{resultsbox}
\textbf{Final Answers:}
\[
\boxed{\matb{\Bb = \Ag^{-1} + \Ag}}
\]
\[
\boxed{\matc{\Cc = -\Ag}}
\]
\end{resultsbox}


% ============================================================================
% PROBLEM 1.7
% ============================================================================

\question{Problem 1.7: Symmetric Matrix Product}\label{quest:1.7}

\begin{problembox}
Show that $AB$ may not be symmetric even though both $A$ and $B$ are symmetric.
\end{problembox}

\begin{stepbox}
\textbf{Counterexample Construction}

We need to find symmetric matrices $\Ag$ and $\Bb$ such that $\Ag\Bb \neq (\Ag\Bb)^T$.

Let:
\[
\Ag = \matg{\begin{bmatrix} 1 & 1 \\ 1 & 0 \end{bmatrix}}, \quad \Bb = \matb{\begin{bmatrix} 0 & 1 \\ 1 & 1 \end{bmatrix}}
\]
\end{stepbox}

\begin{stepbox}
\textbf{Verifying Symmetry of $\Ag$ and $\Bb$}

\[
\Ag^T = \matg{\begin{bmatrix} 1 & 1 \\ 1 & 0 \end{bmatrix}} = \Ag \quad \checkmark
\]

\[
\Bb^T = \matb{\begin{bmatrix} 0 & 1 \\ 1 & 1 \end{bmatrix}} = \Bb \quad \checkmark
\]
\end{stepbox}

\begin{stepbox}
\textbf{Computing the Product}

\[
\Ag\Bb = \matg{\begin{bmatrix} 1 & 1 \\ 1 & 0 \end{bmatrix}} \matb{\begin{bmatrix} 0 & 1 \\ 1 & 1 \end{bmatrix}} = \matc{\begin{bmatrix} 1 & 2 \\ 0 & 1 \end{bmatrix}}
\]

\[
(\Ag\Bb)^T = \matr{\begin{bmatrix} 1 & 0 \\ 2 & 1 \end{bmatrix}}
\]
\end{stepbox}

\begin{stepbox}
\textbf{Comparison}

\[
\Ag\Bb = \matc{\begin{bmatrix} 1 & 2 \\ 0 & 1 \end{bmatrix}} \neq \matr{\begin{bmatrix} 1 & 0 \\ 2 & 1 \end{bmatrix}} = (\Ag\Bb)^T
\]

Therefore, $\Ag\Bb \neq (\Ag\Bb)^T$, demonstrating that $\Ag\Bb$ is \textbf{not symmetric}.
\end{stepbox}

\begin{resultsbox}
\textbf{Conclusion:}

Although both $\Ag$ and $\Bb$ are symmetric matrices, their product $\Ag\Bb$ is \textbf{not symmetric}.

\textbf{Counterexample:}
\[
\Ag = \begin{bmatrix} 1 & 1 \\ 1 & 0 \end{bmatrix}, \quad \Bb = \begin{bmatrix} 0 & 1 \\ 1 & 1 \end{bmatrix}, \quad \Ag\Bb = \begin{bmatrix} 1 & 2 \\ 0 & 1 \end{bmatrix} \neq (\Ag\Bb)^T
\]

\textbf{Note:} The product $\Ag\Bb$ is symmetric if and only if $\Ag\Bb = \Bb\Ag$, i.e., if $\Ag$ and $\Bb$ commute.
\end{resultsbox}


% ============================================================================
% END OF DOCUMENT
% ============================================================================

\end{document}
